\slajdid{01-s010}
\begin{frame}[label=01-s010]
\centering
Funkcija ili preslikavanje iz skupa \(A\) u skup \(B\)\\
je svako pravilo \(f\)\\
po kome se elementu \(x\in A\) pridružuje jedinstveni element \(y\in B\)
\[
f:A\to B \qquad y=f(x)
\]
\begin{columns}[c,onlytextwidth]
\begin{column}{.47\textwidth}
\(A\) – neprazan skup\\
\(n\) – prirodan broj
\end{column}
\begin{column}{.47\textwidth}
\[f:A^n\to A\]
\end{column}
\end{columns}
\raggedright
Operacija – sve funkcije koje preslikavaju \(n\)-ti stepen nekog skupa u sam skup
\begin{itemize}
\item Opis „rečima“
\item Opis funkcionalnim tabelama
\item Opis jednačinama
\item Šematski prikaz korišćenjem simbola operacija - Električna šema
\end{itemize}
\end{frame}
